%% ===========================================================================
%%  Modelling, Control and Simulation of an Inverted Pendulum on a Cart
%%  Dynamics -- Week 1 Course Project
%%  Compile on Overleaf with pdfLaTeX (default). Main file: main.tex
%%  Figures: upload them to the folder "figures/" with the file names given in
%%  each placeholder box. When a file exists it replaces the box automatically.
%% ===========================================================================
\documentclass[10pt,a4paper]{article}

% ---------------- Page layout and fonts ----------------
\usepackage[a4paper,top=1.8cm,bottom=1.8cm,left=1.8cm,right=1.8cm,headsep=0.4cm,footskip=0.8cm]{geometry}
\usepackage[T1]{fontenc}
\usepackage{lmodern}
\usepackage{microtype}

% ---------------- Mathematics ----------------
\usepackage{amsmath,amssymb}

% ---------------- Figures, tables, colours ----------------
\usepackage{graphicx}
\usepackage{float}
\usepackage{array}
\usepackage[table]{xcolor}
\usepackage{caption}
\usepackage{subcaption}
\usepackage{enumitem}

% ---------------- Headings, header/footer, links ----------------
\usepackage{titlesec}
\usepackage{fancyhdr}
\usepackage{lastpage}
\usepackage[hidelinks]{hyperref}

% ======================= EDIT YOUR DETAILS HERE =======================
\newcommand{\studentname}{Mary Sam}
\newcommand{\studentnumber}{}  % student number removed for the public copy
\newcommand{\coursename}{Dynamics}
\newcommand{\reportdate}{October 2026}
% =====================================================================

% ---------------- Colours ----------------
\definecolor{navy}{RGB}{20,45,90}
\definecolor{headgrey}{RGB}{225,231,240}
\definecolor{rulegrey}{RGB}{150,150,150}
\definecolor{boxgrey}{RGB}{245,245,245}

% ---------------- Section style ----------------
\titleformat{\section}{\large\bfseries\color{navy}}{\thesection}{0.6em}{}[\vspace{-2pt}{\color{navy}\titlerule[0.6pt]}]
\titleformat{\subsection}{\normalsize\bfseries\color{navy}}{\thesubsection}{0.5em}{}
\titlespacing*{\section}{0pt}{9pt}{4pt}
\titlespacing*{\subsection}{0pt}{6pt}{2pt}

% ---------------- Header / footer ----------------
\pagestyle{fancy}
\fancyhf{}
\fancyhead[L]{\small\color{navy}Inverted Pendulum on a Cart -- Modelling, PID Control and Simscape Simulation}
\fancyhead[R]{\small\color{navy}\studentname}
\fancyfoot[C]{\small Page \thepage\ of \pageref{LastPage}}
\renewcommand{\headrulewidth}{0.4pt}
\renewcommand{\headrule}{\color{navy}\hrule width\headwidth height\headrulewidth}

% ---------------- Spacing ----------------
\setlength{\parindent}{0pt}
\setlength{\parskip}{3pt}
\setlist{nosep,leftmargin=1.3em}
\setlength{\abovedisplayskip}{4pt}
\setlength{\belowdisplayskip}{4pt}
\setlength{\abovecaptionskip}{3pt}
\setlength{\intextsep}{6pt}
\captionsetup{font=small,labelfont={bf,color=navy},labelsep=period}
\captionsetup[table]{position=top}
\renewcommand{\arraystretch}{1.18}
\arrayrulecolor{rulegrey}

% ---------------- Table header row helper ----------------
\newcommand{\hdr}{\rowcolor{headgrey}}

% ---------------- Figure placeholder ----------------
% \figslot{file name}{width}{height}{what the figure should show}
% If figures/<file name> exists it is inserted; otherwise a grey box describes it.
\newcommand{\figslot}[4]{%
  \IfFileExists{figures/#1}{%
    \includegraphics[width=#2,height=#3,keepaspectratio]{figures/#1}%
  }{%
    \setlength{\fboxsep}{0pt}\color{rulegrey}%
    \fbox{\colorbox{boxgrey}{\parbox[c][#3][c]{#2}{\centering\color{black}\small
      \textbf{Insert figure here}\\[2pt]\texttt{figures/\detokenize{#1}}\\[4pt]
      \begin{minipage}{0.88\linewidth}\centering\footnotesize\itshape #4\end{minipage}}}}%
  }%
}

\begin{document}

%% =====================================================================
%%  TITLE BLOCK
%% =====================================================================
\thispagestyle{fancy}
\begin{center}
{\Large\bfseries\color{navy}\MakeUppercase{\coursename}}\enspace{\large\color{navy}--\enspace Course Project}\\[4pt]
{\LARGE\bfseries\color{navy} Modelling, Control and Simulation of an\\[2pt] Inverted Pendulum on a Cart}\\[6pt]
{\normalsize\textbf{\studentname} \quad$|$\quad \reportdate}
\end{center}
\vspace{-2pt}

%% =====================================================================
\section{Introduction and objectives}
%% =====================================================================
The inverted pendulum on a cart is the classic benchmark for unstable, under-actuated systems; the same
balancing problem appears in two-wheeled robots, self-balancing vehicles and rockets during launch.
In this project a PID controller acting on the cart force keeps a 6.64\,kg steel pendulum upright on a
75.3\,kg cart. The objectives were to:
\begin{enumerate}
  \item derive the nonlinear equations of motion with the Lagrangian method and linearise them about the
        upright equilibrium;
  \item design a PID controller for the linearised model;
  \item simulate the closed loop in MATLAB and Simulink;
  \item build the system in SolidWorks, simulate it in Simscape Multibody with the same gains, and compare
        the CAD-based response with the analytical one.
\end{enumerate}

%% =====================================================================
\section{System and CAD model}
%% =====================================================================
The system (Fig.~\ref{fig:cad}) consists of a fixed platform with two rails, a steel cart on four wheels,
and a pendulum made of a rod, a 100\,mm steel sphere and a pivot pin, hinged on top of the cart. All
parameters were taken from the SolidWorks mass properties (Table~\ref{tab:params}). The platform is fixed
to the ground, so its mass does not enter the dynamics and is excluded. In the assembly the pendulum is
not exactly vertical: it leans $\theta_0 = 0.02385$\,rad ($1.37^\circ$) towards $+x$. This tilt of the CAD pose
is the \textbf{initial condition} of every simulation: the pendulum is released from $1.37^\circ$ with the cart
at rest, and the controller must bring it back to upright ($\theta = 0$).

\begin{figure}[H]
\centering
\figslot{fig1_cad_assembly.png}{8cm}{4.4cm}{SolidWorks assembly (or Simscape Mechanics Explorer view) with
the cart position $x$, the pendulum angle $\theta$ and the force $F$ labelled.}
\caption{SolidWorks assembly of the inverted pendulum on a cart: rail platform, cart on four wheels and
pendulum (rod and sphere) hinged on the cart. $x$ is the cart position along the rail, $\theta$ the pendulum
tilt from vertical and $F$ the control force on the cart. In this pose the pendulum leans $\theta_0 = 1.37^\circ$,
which is used as the initial tilt in all simulations.}
\label{fig:cad}
\end{figure}

\begin{table}[H]
\centering\small
\caption{Model parameters obtained from the CAD model (plain carbon steel, $\rho = 7800$\,kg/m$^3$).}
\label{tab:params}
\setlength{\tabcolsep}{4.5pt}
\begin{tabular}{|c|l|r|l|}
\hline \hdr
\textbf{Symbol} & \textbf{Quantity} & \textbf{Value} & \textbf{Source} \\ \hline
$m$ & Pendulum mass (rod + sphere + pin) & 6.640 kg & CAD \\ \hline
$l$ & Pivot to pendulum centre of mass & 0.36605 m & CAD \\ \hline
$I_{cm}$ & Pendulum inertia about its centre of mass & 0.25241 kg\,m$^2$ & CAD \\ \hline
$J$ & Pendulum inertia about the pivot ($I_{cm}+ml^2$) & 1.14218 kg\,m$^2$ & Parallel axis \\ \hline
$m_c$ & Cart body with axles & 71.453 kg & CAD \\ \hline
$m_w$ / $I_w$ / $R_w$ & Wheel mass / inertia / radius & 0.588 kg / $5.88{\times}10^{-4}$ kg\,m$^2$ / 0.04 m & CAD \\ \hline
$M$ & Effective cart mass (rolling wheels, Sec.~3) & 75.275 kg & Calculated \\ \hline
$g$ & Gravity & 9.81 m/s$^2$ & -- \\ \hline
$b,\ c$ & Rail and pivot friction & 0 & Assumption \\ \hline
$\theta_0$ & Initial tilt (CAD pose) & 0.02385 rad (1.37$^\circ$) & CAD \\ \hline
\end{tabular}
\end{table}

%% =====================================================================
\section{Analytical modelling (Phase 1)}
%% =====================================================================
\textbf{Assumptions.} The bodies are rigid with the CAD mass properties; motion is planar; there is no
friction in the rail or pivot; the wheels roll without slip; the actuator applies the force directly to
the cart and the angle is measured exactly (no delay or noise).

\textbf{Coordinates and energies.} The generalised coordinates are the cart position $x$ and the pendulum
angle $\theta$ from the upright. The pendulum centre of mass is at $(x+l\sin\theta,\ l\cos\theta)$, so
\begin{equation}
T = \tfrac12 (M+m)\dot x^2 + m l\,\dot x\,\dot\theta\cos\theta + \tfrac12 J\dot\theta^2,\qquad
V = m g l\cos\theta,\qquad \mathcal L = T - V .
\end{equation}

\textbf{Lagrange's equations.} Applying
$\dfrac{d}{dt}\Big(\dfrac{\partial\mathcal L}{\partial\dot q_i}\Big)-\dfrac{\partial\mathcal L}{\partial q_i}=Q_i$
with $Q_x = F$ and $Q_\theta = 0$ gives the coupled nonlinear equations of motion
\begin{align}
(M+m)\ddot x + m l\cos\theta\,\ddot\theta - m l\,\dot\theta^2\sin\theta &= F, \label{eq:eom1}\\
m l\cos\theta\,\ddot x + J\ddot\theta - m g l\sin\theta &= 0 . \label{eq:eom2}
\end{align}

\textbf{Linearisation.} About the upright equilibrium ($\sin\theta\approx\theta$, $\cos\theta\approx1$,
$\dot\theta^2\approx0$), equations (\ref{eq:eom1})--(\ref{eq:eom2}) become
$(M+m)\ddot x + ml\ddot\theta = F$ and $ml\ddot x + J\ddot\theta - mgl\,\theta = 0$. With state
$\mathbf s=[x,\ \dot x,\ \theta,\ \dot\theta]^\mathsf T$, input $F$, output $\mathbf y=[x,\ \theta]^\mathsf T$ and
$\Delta=(M+m)J-(ml)^2 = 87.654$\,kg$^2$m$^2$:
\begin{equation}
A=\begin{bmatrix} 0&1&0&0\\ 0&0&-\frac{(ml)^2g}{\Delta}&0\\ 0&0&0&1\\ 0&0&\frac{(M+m)mgl}{\Delta}&0\end{bmatrix}
=\begin{bmatrix} 0&1&0&0\\ 0&0&-0.6613&0\\ 0&0&0&1\\ 0&0&22.284&0\end{bmatrix},\qquad
B=\begin{bmatrix}0\\ \frac{J}{\Delta}\\ 0\\ -\frac{ml}{\Delta}\end{bmatrix}=\begin{bmatrix}0\\ 0.013031\\ 0\\ -0.027731\end{bmatrix},
\end{equation}
\begin{equation}
\dot{\mathbf s} = A\mathbf s + BF,\qquad \mathbf y = C\mathbf s + DF,\qquad
C=\begin{bmatrix}1&0&0&0\\0&0&1&0\end{bmatrix},\quad D=\begin{bmatrix}0\\0\end{bmatrix}.
\end{equation}
The pendulum-angle transfer function is
\begin{equation}
G_\theta(s)=\frac{\Theta(s)}{F(s)} = \frac{-0.02773}{s^2-22.28},\qquad \text{poles } s=\pm4.72\ \text{rad/s}
\ (\text{plus a double pole at } 0 \text{ for the cart}).
\end{equation}
The pole at $+4.72$\,rad/s makes the upright position unstable, so feedback control is required. The
derivation was verified with the Symbolic Math Toolbox (symbolic and numerical $A$, $B$ agree to
$4\times10^{-15}$).

\textbf{Check of the CAD quantities.}
\begin{itemize}
  \item $l = 0.36605$\,m (pivot at 0.5000\,m, centre of mass at 0.1339\,m along the rod axis). Hand check:
        a 100\,mm steel sphere (4.084\,kg) at 451.6\,mm from the pivot plus the rest of the mass as a uniform
        rod gives $l = 0.3647$\,m, within 0.4\,\%.
  \item $J = I_{cm} + ml^2 = 0.25241 + 6.6404\times0.36605^2 = 1.14218$\,kg\,m$^2$. Shifting $I_{cm}$ to the
        part origin with the same theorem gives 0.371554\,kg\,m$^2$, matching SolidWorks' output-frame value
        (0.371553\,kg\,m$^2$).
  \item $M = 71.453 + 4(0.5881 + 5.881\times10^{-4}/0.04^2) = 75.275$\,kg: rolling without slip adds
        $I_w/R_w^2 = 0.368$\,kg per wheel.
\end{itemize}

%% =====================================================================
\section{Control design (Phase 2)}
%% =====================================================================
A parallel PID controller with a filtered derivative acts on the angle error $e = 0-\theta$:
\begin{equation}
F(s) = \Big(K_p + \frac{K_i}{s} + \frac{K_d\,s}{T_f s+1}\Big)E(s),\qquad T_f = 0.01\ \text{s}\ (N = 100\ \text{rad/s}).
\end{equation}
The gains are negative because $G_\theta(s)$ has a negative gain: pushing the cart forward tips the
pendulum backward. The gains were tuned to minimise the ITAE of $\theta$ subject to a phase margin of at
least $60^\circ$, a peak sensitivity $M_s \le 1.2$ and a peak force of at most 100\,N for the CAD tilt
(Table~\ref{tab:pid}). The main script also runs MATLAB's PID Tuner algorithm (\texttt{pidtune}) on the same
plant at the same crossover frequency and phase margin as a cross-check.

\begin{table}[H]
\centering\small
\caption{PID gains and loop properties.}
\label{tab:pid}
\begin{tabular}{|l|c||l|c|}
\hline \hdr
\textbf{Parameter} & \textbf{Value} & \textbf{Property} & \textbf{Value} \\ \hline
Proportional gain $K_p$ & $-4192.9$ & Phase margin & $60.0^\circ$ at 17.95 rad/s \\ \hline
Integral gain $K_i$ & $-8365.6$ & Peak sensitivity $M_s$ & 1.16 \\ \hline
Derivative gain $K_d$ & $-645.6$ & Gain margin (gain reduction) & 0.30 ($-10.5$ dB) \\ \hline
Filter time constant $T_f$ & 0.01 s & Robustness & Stable for $\pm20\,\%$ in $m$ and $M$ \\ \hline
\end{tabular}
\end{table}

%% =====================================================================
\section{MATLAB and Simulink results}
%% =====================================================================
The closed loop was simulated from the CAD tilt $\theta_0 = 1.37^\circ$ with the cart at rest, using the
nonlinear equations (\ref{eq:eom1})--(\ref{eq:eom2}) and the linearised model, and with the Simulink model
\texttt{pendulum\_pid.slx} (State-Space plant and PID Controller block, Fig.~\ref{fig:phase2}a).

\begin{figure}[H]
\centering
\begin{subfigure}[b]{0.53\linewidth}\centering
  \figslot{fig2a_simulink_model.png}{\linewidth}{5.6cm}{Simulink block diagram \texttt{pendulum\_pid.slx}:
  ThetaRef, Sum, PID, State-Space plant, Demux and scopes.}
  \caption{Simulink model with the PID control loop}
\end{subfigure}\hfill
\begin{subfigure}[b]{0.43\linewidth}\centering
  \figslot{fig2b_closed_loop.pdf}{\linewidth}{5.8cm}{Three stacked panels from
  \texttt{main\_pid\_inverted\_pendulum.m}: (a) $\theta(t)$ nonlinear and linearised, (b) cart position $x(t)$,
  (c) control force $F(t)$, 0--8 s.}
  \caption{Closed-loop response (analytical model)}
\end{subfigure}
\caption{Phase 2: (a) Simulink model \texttt{pendulum\_pid.slx}; (b) pendulum angle, cart position and
control force of the analytical model. The pendulum starts from the initial tilt $\theta_0 = 1.37^\circ$ of the
CAD pose with the cart at rest; the nonlinear and linearised curves coincide.}
\label{fig:phase2}
\end{figure}

\begin{table}[H]
\centering\small
\caption{Closed-loop performance of the analytical model.}
\label{tab:phase2}
\begin{tabular}{|l|c|l|}
\hline \hdr
\textbf{Quantity} & \textbf{Value} & \textbf{Interpretation} \\ \hline
Settling time of $\theta$ (2\,\%) & 1.22 s & Pendulum upright after 1.22 s \\ \hline
Undershoot past upright & 27.4\,\% ($-0.375^\circ$ at 0.54 s) & One small swing, no oscillation \\ \hline
Peak control force & 100.1 N & Largest push, at $t=0$ \\ \hline
Cart peak / final position & 31.1 mm / 29.4 mm & Cart moves under the pendulum, then stops \\ \hline
Closed-loop poles & $-77.1$, $-16.7$, $-3.08\pm2.92j$ & All in the left half-plane: stable \\ \hline
\end{tabular}
\end{table}

The pendulum balances: $\theta$ crosses zero at 0.28\,s, swings $0.375^\circ$ past upright and stays within
2\,\% of $\theta_0$ after 1.22\,s. The cart comes to rest 29.4\,mm from its start once the pendulum is upright.
The nonlinear and linearised angles differ by less than $2\times10^{-6}$\,rad, and the Simulink model gives
the same result.

%% =====================================================================
\section{Simscape model (Phase 3)}
%% =====================================================================
The SolidWorks assembly was exported with the Simscape Multibody Link plug-in and imported with
\texttt{smimport} (Fig.~\ref{fig:simscape}). The platform is fixed to the world frame; the cart moves on a
\textbf{prismatic joint} along the rail; the pendulum turns on a \textbf{revolute joint} at the pivot; the
wheels turn freely on their axles. All joint damping is zero. The imported hinge reads $90^\circ$ when the
pendulum is upright, so an offset of $-\pi/2$ converts it to $\theta$. The same PID Controller block and gains
as in Phase 2 drive the prismatic joint, and the simulation starts from the CAD pose ($\theta_0 = 1.37^\circ$).

\begin{figure}[H]
\centering
\figslot{fig3_simscape_model.png}{\linewidth}{7.2cm}{Simscape block diagram \texttt{PendulumAssy.slx}:
imported CAD bodies, prismatic (cart) and revolute (pendulum) joints, sensors, $-\pi/2$ Bias block and the
PID Controller loop.}
\caption{Phase 3: (a, left) CAD-based Simscape Multibody model \texttt{PendulumAssy.slx} with the PID
control loop; (b, right) its scope output: pendulum angle $\theta$ [rad], cart position $x$ [m] and control
force $F$ [N] over 8\,s, starting from the same initial tilt $\theta_0 = 0.0239$\,rad ($1.37^\circ$).}
\label{fig:simscape}
\end{figure}

\begin{figure}[H]
\centering
\figslot{fig4_comparison.png}{0.62\linewidth}{8.4cm}{Output of \texttt{compare\_pid\_vs\_assy.m}:
$\theta(t)$, $x(t)$ and $F(t)$ for Simscape (solid) and Simulink (dashed) on the same axes, plus the
angle difference, 0--4 s.}
\caption{Comparison of the Simulink (analytical) and Simscape (CAD) results with identical PID gains:
(a) pendulum angle, starting from the initial tilt $\theta_0 = 1.37^\circ$, (b) cart position, (c) control
force, (d) angle difference $\theta_{\mathrm{Simscape}}-\theta_{\mathrm{Simulink}}$.}
\label{fig:cmp}
\end{figure}

\begin{table}[H]
\centering\small
\caption{Analytical (Simulink) and CAD (Simscape) results.}
\label{tab:cmp}
\begin{tabular}{|l|c|c|c|c|}
\hline \hdr
\textbf{Metric} & \textbf{Simulink} & \textbf{Simscape} & \textbf{Difference} & \textbf{Difference (\%)} \\ \hline
Settling time of $\theta$, 2\,\% [s] & 1.219 & 1.224 & $+0.005$ & $+0.41$ \\ \hline
Undershoot [\% of $\theta_0$] & 27.43 & 27.05 & $-0.38$ & $-1.37$ \\ \hline
Peak force [N] & 100.12 & 99.97 & $-0.15$ & $-0.15$ \\ \hline
Cart peak / final position [mm] & 31.05 / 29.26 & 30.98 / 29.41 & $-0.07$ / $+0.15$ & $-0.21$ / $+0.51$ \\ \hline
\end{tabular}
\end{table}

%% =====================================================================
\section{Discussion}
%% =====================================================================
\textbf{Convergence and stability.} The open-loop model has an unstable pole at $+4.72$\,rad/s. With the
PID controller all closed-loop poles lie in the left half-plane, the phase margin is $60^\circ$ and $M_s =
1.16$, so the loop is stable and robust: it stays stable for $\pm20\,\%$ changes in either mass. In both
models $\theta$ converges to zero in about 1.22\,s after a single undershoot, and the cart comes to rest.

\textbf{Differences between the analytical and CAD results.} The Simscape curve lies on top of the
Simulink curve (Fig.~\ref{fig:cmp}). The largest angle difference is $0.0067^\circ$ (RMS $0.0014^\circ$),
the settling time differs by 5\,ms and every metric in Table~\ref{tab:cmp} agrees to within 1.4\,\%.

\textbf{Sources of discrepancy.} (i) \emph{Wheel rolling inertia}: the analytical cart mass includes it
($M = 75.28$\,kg), while the Simscape wheels spin freely without rail contact (73.80\,kg); this is the most
likely cause of the 0.38-point difference in undershoot. (ii) \emph{Inertia details}: Simscape uses the full 3-D
inertia tensors from CAD, the analytical model only the planar values. (iii) \emph{Gravity constant}:
9.81 against 9.80665\,m/s$^2$. (iv) \emph{Linearisation}: the Simulink plant is linear; at $1.37^\circ$ this
contributes only about $10^{-6}$\,rad. (v) \emph{Solver tolerance}: ode45 against a variable-step Simscape
solver (relative tolerance $10^{-6}$). (vi) \emph{Zero friction}: both models neglect friction, so their
agreement validates the equations and CAD parameters but not a physical rig.

\textbf{Limitations.} The controller uses only $\theta$, so the cart position is not regulated: after a
disturbance the pendulum recovers within about 1\,s but the cart keeps drifting. The heavy steel cart also
needs about 100\,N for a $1.4^\circ$ tilt. Cart-position feedback (cascade PID or LQR) and a lighter cart would
address both.

%% =====================================================================
\section{Conclusions}
%% =====================================================================
The Lagrangian model, built from CAD-derived parameters, has an unstable pole at $+4.72$\,rad/s. A PID
controller ($K_p=-4192.9$, $K_i=-8365.6$, $K_d=-645.6$) with a $60^\circ$ phase margin balances the pendulum
from a $1.37^\circ$ tilt in 1.22\,s with a peak force of 100\,N, and the cart settles 29\,mm from its start.
The CAD-based Simscape model, with the same gains, agrees with the analytical model to within
$0.0067^\circ$ in angle and 1.4\,\% in all metrics.

{\small\textbf{References.}
[1] MathWorks, \emph{Control System Toolbox: pidtune / PID Tuner}, MATLAB R2025b.
[2] MathWorks, \emph{Simscape Multibody Link and smimport}, MATLAB R2025b.
[3] K.~Ogata, \emph{Modern Control Engineering}, 5th ed., Prentice Hall, 2010.}

%% =====================================================================
%%  PAGE 5: PROJECT FILES AND HOW TO RUN THEM
%% =====================================================================
\clearpage
\section*{Appendix: project files and how to reproduce the results}
\addcontentsline{toc}{section}{Appendix: project files}
All files are in the submitted folder. They need \textbf{MATLAB R2025b} with Simulink, Control System
Toolbox, Symbolic Math Toolbox, Simscape and Simscape Multibody. SolidWorks 2026 is needed only to open or
re-export the CAD model.

\begin{table}[H]
\centering\small
\caption{Project files.}
\begin{tabular}{|>{\raggedright}p{0.30\linewidth}|>{\raggedright}p{0.10\linewidth}|>{\raggedright\arraybackslash}p{0.52\linewidth}|}
\hline \hdr
\textbf{File} & \textbf{Phase} & \textbf{Purpose} \\ \hline
\multicolumn{3}{|l|}{\cellcolor{boxgrey}\textit{MATLAB scripts}} \\ \hline
\texttt{params.m} & All & All model parameters from the CAD model, and the linearised $A$, $B$, $C$, $D$ matrices. Called by the other scripts. \\ \hline
\texttt{derive\_symbolic.m} & 1 & Lagrangian derivation, nonlinear equations of motion and linearisation with the Symbolic Math Toolbox; checks the result against \texttt{params.m}. \\ \hline
\texttt{main\_pid\_inverted\_pendulum.m} & 2 & Transfer function, PID gains, loop margins, nonlinear and linear simulations and the 3-panel figure. Saves \texttt{pid\_gains.mat} and \texttt{analytical\_results.mat}. \\ \hline
\texttt{optimize\_pid.m} & 2 & Optional: reproduces the PID gains (ITAE minimisation with margin and force constraints). Takes a few minutes. \\ \hline
\texttt{build\_simulink\_pid\_model.m} & 2 & Builds and runs the Simulink model \texttt{pendulum\_pid.slx}. \\ \hline
\texttt{compare\_pid\_vs\_assy.m} & 3 & Runs \texttt{pendulum\_pid.slx} and \texttt{PendulumAssy.slx}, plots both on the same axes and prints the comparison table. \\ \hline
\multicolumn{3}{|l|}{\cellcolor{boxgrey}\textit{Models and data}} \\ \hline
\texttt{pendulum\_pid.slx} & 2 & Simulink model: linearised State-Space plant with the PID Controller block. \\ \hline
\texttt{PendulumAssy.slx} & 3 & Simscape Multibody model imported from SolidWorks, with prismatic and revolute joints and the PID loop. \\ \hline
\texttt{PendulumAssyDisturb.slx} & 3 & Same model with a torque pulse on the pendulum (disturbance test). \\ \hline
\texttt{PendulumAssy\_DataFile.m} & 3 & Data written by \texttt{smimport} (joint frames, masses, inertias); needed by both Simscape models. \\ \hline
\texttt{PendulumAssy.xml}, \texttt{*\_sldprt.STEP} & 3 & SolidWorks export (assembly structure and 3-D geometry of each part). \\ \hline
\texttt{pid\_gains.mat}, \texttt{analytical\_results.mat} & 2 & Saved gains and analytical simulation results (created by the main script). \\ \hline
\multicolumn{3}{|l|}{\cellcolor{boxgrey}\textit{CAD files}} \\ \hline
\texttt{Pendulum-Assy.SLDASM} & 3 & SolidWorks assembly (open together with the part files). \\ \hline
\texttt{Platform}, \texttt{Wheels}, \texttt{Wheel 1}, \texttt{InvertedPendulum} \texttt{.SLDPRT} & 3 & SolidWorks parts: rail platform, cart body, wheel and pendulum. \\ \hline
\end{tabular}
\end{table}

\subsection*{How to run}
\begin{enumerate}
  \item Copy the whole folder to your computer, open MATLAB and set the folder as the \emph{Current Folder}.
  \item \textbf{Phase 1:} run \texttt{derive\_symbolic}. The command window shows the nonlinear equations, the
        $A$ and $B$ matrices and two norms close to $10^{-15}$.
  \item \textbf{Phase 2:} run \texttt{main\_pid\_inverted\_pendulum}. It prints the open-loop poles, the
        transfer function, the gains, the phase margin and the performance figures, and opens the 3-panel
        response plot (Fig.~\ref{fig:phase2}b).
  \item Run \texttt{build\_simulink\_pid\_model} (or open \texttt{pendulum\_pid.slx} and press \emph{Run}) to
        simulate the Simulink model; open the scopes to see $\theta$, $x$ and $F$.
  \item \textbf{Phase 3:} run \texttt{PendulumAssy\_DataFile}, then open \texttt{PendulumAssy.slx} and press
        \emph{Run}. Mechanics Explorer shows the 3-D animation of the pendulum balancing; the scope shows
        $\theta$, $x$ and $F$.
  \item \textbf{Comparison:} run \texttt{compare\_pid\_vs\_assy}. It produces Fig.~\ref{fig:cmp} (saved as
        \texttt{fig\_simulink\_vs\_simscape.png}) and the values of Table~\ref{tab:cmp} (saved as
        \texttt{simulink\_vs\_simscape.csv}).
  \item \textbf{Optional:} run \texttt{optimize\_pid} to re-derive the gains, or open
        \texttt{PendulumAssyDisturb.slx} and press \emph{Run} for the disturbance test.
\end{enumerate}

\end{document}
